%% ===========================================================================
%%  6. Conclusion
%% ===========================================================================

\section{Conclusion}
\label{sec:conclusion}

We presented LP-Cos-CP, a framework that replaces opaque single-genre
predictions with interpretable, user-controllable prediction sets. On GTZAN and
FMA Small, our method achieves empirical coverage tightly matching nominal
targets while reducing set sizes by up to 55\% compared to centroid
baselines. The $\alpha$ parameter functions as a principled trust slider,
enabling users to configure the reliability--informativeness trade-off for their
specific task.

Our comparison of calibration strategies reveals that distribution-free
calibration is not merely a statistical nicety but a practical requirement for
predictable user-facing behavior: parametric assumptions (Gaussian,
Bootstrap) produce interface behavior that is either uninformative or
untrustworthy. CP's empirical $p$-value calibration remains robust across
datasets and score distributions.

The prediction set provides a template for bridging the gap between ML
performance and human usability. Future work includes user studies evaluating
prediction sets in interactive music tools, novelty detection for
out-of-distribution tracks, and prototype personalization through user feedback.
We believe this interface paradigm extends naturally to other perceptual
classification domains.
